\documentclass[11pt,a4paper]{article}

\usepackage[utf8]{inputenc}
\usepackage[T1]{fontenc}
\usepackage[french]{babel}
\usepackage[margin=2.5cm]{geometry}
\usepackage{enumitem}
\usepackage{hyperref}

\title{Résumé du projet \texttt{sentiment-app}}
\author{Hiba Ayatallah et Manal Wassila}
\date{7 octobre 2026}

\begin{document}
\maketitle

\section{Présentation}

\texttt{sentiment-app} est le dépôt de travail du binôme pour le cours
\emph{Développement et déploiement d'applications intelligentes}. Le projet
porte sur une application d'analyse de sentiments à partir d'avis. La première
semaine a permis de mettre en place le dépôt Git, l'environnement Python, les
tests et le cadrage du projet. Les prochaines étapes prévues comprennent une API
FastAPI, une interface Streamlit et la gestion des données et modèles.

\section{Travail réalisé}

Le dépôt comprend les consignes et guides du TP, les fiches de cadrage, un jeu
de données d'exemple, des scripts de diagnostic et une suite de tests avec
\texttt{pytest}. Les tests valident notamment l'environnement, les fiches de
cadrage et l'historique Git du binôme.

\section{Difficultés rencontrées}

\begin{enumerate}[leftmargin=*]
    \item \textbf{Pytest absent du Python global.} La première commande
    \texttt{python -m pytest} a utilisé l'interpréteur système, dans lequel le
    module \texttt{pytest} n'était pas installé. La solution a été d'utiliser
    directement l'interpréteur de l'environnement virtuel du projet :
    \texttt{\detokenize{.\.venv\Scripts\python.exe -m pytest -q}}.

    \item \textbf{Commit de la seconde membre absent de l'historique local.}
    Le test complet signalait qu'un seul auteur de commit était visible,
    alors que le commit de Hiba avait déjà été envoyé sur GitHub. Après un
    \texttt{git fetch origin} et une mise à jour en fast-forward avec
    \texttt{git pull --ff-only origin main}, le dépôt local a récupéré le
    commit distant et le test d'historique a réussi.
\end{enumerate}

\section{Résultat}

Après la synchronisation avec GitHub, la suite complète s'est terminée avec
\textbf{38 tests réussis}. Les 12 tests de cadrage avaient également réussi.
Le dépôt local était ensuite synchronisé avec la branche \texttt{main} de
GitHub.

\end{document}
